\section{Construction of a stationary equality}
\label{sec:stationary}

We work at the differentiability point supplied by
Lemma~\ref{lem:critical-point}, and use the exact, positive-deficit equality
from Lemma~\ref{lem:positive-deficit}. Thus the upper exponent and the
exponent in the plank parameter are both $h$. Throughout this section,
\[
 x=d-1-h\ge0,\qquad y=1+h,\qquad x+y=d,
 \qquad y-x=q+2h>0.
\]
The word equality refers to a realizing sequence in the normalized
logarithmic sense of Section~\ref{sec:critical}. All local masses retain the
original index weights. In particular, taking an extensive regular refinement does not
renormalize those weights.

The two rescalings from Section~\ref{sec:local} retain complementary
parts of a temporal profile. An isotropic cut keeps the scales above a
chosen block, with full time grafted below them; an anisotropic cut keeps
the scales within that block and removes an initial flat interval. At
exact criticality, saturated packets make both rescalings attain the
sharp inequality. We call the resulting models outer and inner equalities.

Parameter derivatives first select the packet shapes, and a packing
argument makes their normals coherent enough for the inner rescaling.
We then compare the delays and deficits of the two resulting profiles.
This forces a profile that is flat and then linear, and whose normalized
outer models reproduce the same profile. That self-reproduction, together
with a constant rate of change of packet widths, is the stationarity
constructed here.

\subsection{Derivative capture and outer equality models}

At logarithmic scale $\kappa$, put $K=N^\kappa$ and write the local bound from
Proposition~\ref{prop:local-bound} as
\begin{equation}
 T(\kappa,a,b)
 =mN^{-\Pi(\kappa)}K^{-h}a^xb^y,
 \qquad 1\le a\le b\le K.
 \label{eq:stationary-local-threshold}
\end{equation}
A saturated assignment at this scale is a disjoint assignment of active
indices to planks on each $K$-block whose occupied packets have mass
$\eqexp T(\kappa,a,b)$, with one common width type. Disjointness concerns
the assigned indices on a fixed temporal block; geometric planks need not
be disjoint.

Here and below a prefix means the logarithmic scale interval
$[0,\kappa]$, realized on every occupied physical $K$-block; it does not
mean only the physical interval starting at time zero. The next lemma
uses changes of the critical parameters to select saturated shapes. The
$p$ derivative measures temporal deficit, while the $z$ derivative also
detects the thin spatial width.

\begin{lemma}[Derivative capture]
\label{lem:derivative-capture}
Every equality profile has deficit $\alpha\le v=-h_p$. In particular,
$v>0$. Fix an equality, a prefix $\kappa>0$, and a perturbation direction
$(\dot p,\dot z)$. On an extensive refinement there are saturated
assignments with limiting shape parameters
\[
 \mathfrak a=\log_N a,\qquad
 \mathfrak u=\log_N(K/b),\qquad
 \chi=\mathfrak a+\mathfrak u,
\]
satisfying
\begin{equation}
 \bigl(F(\kappa)-v\chi\bigr)\dot p
 +\left(h_z\chi+
   \int_0^\kappa\mathcal D_q(q,F'(s))\,ds
   -\mathfrak a\right)\dot z\ge0.
 \label{eq:derivative-capture}
\end{equation}
Consequently there are selections with $\chi\le F(\kappa)/v$ and
selections with $\chi\ge F(\kappa)/v$. These assertions remain valid
after any previously imposed extensive regular refinement.
\end{lemma}

\begin{proof}
Decrease $p$ by a positive amount $\varepsilon$. Differentiability gives
$h(p-\varepsilon,z)=h+v\varepsilon+o(\varepsilon)$, with $v\ge0$.
Increasing the clustering exponent cannot increase $\Delta_h$: for each
plank, its dependence on that exponent is through $(a/b)^h$. The cost
changes by $-\alpha\varepsilon$. Applying the perturbed upper bound to an
equality therefore gives $\alpha\le v$. Since a positive-deficit equality
exists, $v>0$.

For the local assertion fix the perturbation size first. Denote the
nearby parameters by primes, and let $\Pi'$ denote the cost with the
changed integrand and the same profile. On each prefix block greedily
assign and remove the remaining active indices in planks whose mass is
at least
\begin{equation}
 mN^{-\mathrm{tol}-\Pi'(\kappa)}K^{-h'}
       a^{d'-1-h'}b^{1+h'}.
 \label{eq:perturbed-capture-threshold}
\end{equation}
Here $\mathrm{tol}>0$ is fixed during this operation. If the captured
events were not extensive, the remainder would be extensive. Regularize
that remainder to its inherited profile and apply the perturbed critical
bound separately on the prefix blocks and bounded spatial charts. Its
plank parameter is bounded by the threshold used in the greedy removal,
so this bound gives aggregate raw multiplicity at most
$mN^{-\mathrm{tol}+o(1)}$. The extensive remainder has aggregate raw
multiplicity at least $mN^{-o(1)}$, a contradiction. Spatial charts may be
assigned by position at the block center; bounded slopes give bounded
overlap in their support counts.

The captured graph is therefore extensive. Pigeonhole a common dyadic
width type and use Lemma~\ref{lem:regularization} to preserve the packet
lower counts. Comparing \eqref{eq:perturbed-capture-threshold} with
\eqref{eq:stationary-local-threshold} yields
\begin{equation}
 \Pi'(\kappa)-\Pi(\kappa)
 +(h'-h)\chi-(d'-d)\mathfrak a\ge-\mathrm{tol}.
 \label{eq:finite-parameter-capture}
\end{equation}
The first term records the change in temporal cost. The remaining terms
record the resolution lost in the cut and the change of its thin-width
exponent. The shapes lie in the compact triangle
$\mathfrak a,\mathfrak u\ge0$, $\mathfrak a+\mathfrak u\le\kappa$.
Take the perturbation to zero along the specified direction, with
$\mathrm{tol}$ of smaller order, and pass to a convergent sequence of
shapes. First choosing the resolution sufficiently large at each fixed
perturbation makes the preceding exponent errors smaller than
$\mathrm{tol}$. The captured masses are squeezed to saturation in the
limit.

Since $z=-q$ and $d=2+z$,
\[
 \partial_p\Pi(\kappa)=F(\kappa),\qquad
 \partial_z\Pi(\kappa)
   =\int_0^\kappa\mathcal D_q(q,F'(s))\,ds,
 \qquad d_z=1.
\]
The averaging in $\mathcal D$ makes its $q$ derivative available at the
interior parameter in use. Differentiating
\eqref{eq:finite-parameter-capture} proves
\eqref{eq:derivative-capture}. Taking $(\dot p,\dot z)=(1,0)$ and
$(-1,0)$ gives the two final selections. The same argument applies to a
preceding extensive regular refinement because its profile, equality
exponents, and inherited local upper bounds are unchanged.
\end{proof}

We record explicitly the equality obtained from a saturation. This also
fixes the scale bookkeeping for the iteration. If $\chi<1$, the isotropic
parent construction in Proposition~\ref{prop:local-bound} has resolution
\[
 N_{\mathrm{out}}=\frac{Nb}{Ka}=N^{1-\chi}.
\]
Its physical cell width in the old cell units is $a=N^{\mathfrak a}$.
The raw multiplicity supplied by a saturated packet is
\begin{equation}
 T(\kappa,a,b)\frac ab
 =mN^{-\Pi(\kappa)-h\chi+d\mathfrak a}
 \leexp m_{\mathrm{out}}.
 \label{eq:outer-multiplicity}
\end{equation}
Pulling back an arbitrary full-time plank multiplies both widths in cell
units by $a$, so its parameter is at most $\Delta_h a^d$. Filling the
selected $K$-blocks drops exactly the cost below $\kappa$. Thus the
critical upper estimate is
\[
 \Delta_h a^dN_{\mathrm{out}}^h
       N^{\Pi(1)-\Pi(\kappa)}
 =mN^{-\Pi(\kappa)-h\chi+d\mathfrak a}.
\]
Consequently the outer model is an equality. In normalized log
coordinates its profile is
\begin{equation}
 F_{\mathrm{out}}(r)=
 \begin{cases}
 0,&0\le r\le(\kappa-\chi)/(1-\chi),\\[2mm]
 \displaystyle
 \frac{F(\chi+(1-\chi)r)-F(\kappa)}{1-\chi},
    &(\kappa-\chi)/(1-\chi)\le r\le1.
 \end{cases}
 \label{eq:outer-profile}
\end{equation}
In particular its deficit is $(\alpha-F(\kappa))/(1-\chi)$.
The full-time isotropic parents have one owner per index. Later local
plank bounds pull back with the same multiplication of widths, and their
active indices belong to the corresponding physical blocks.

\subsection{Finite chains and orientation control}

Call an equality nontrivial if its total deficit $F(1)$ is positive.
By Lemma~\ref{lem:flat-delay}, such an equality has an initial flat
interval whose endpoint $\tau$ satisfies $0<\tau<1$.

\begin{lemma}[Compatible finite cuts]
\label{lem:nested-cuts}
Fix a nontrivial equality and a finite increasing mesh
$\tau=\kappa_0<\kappa_1<\cdots<\kappa_J<1$. There is an extensive
refinement with simultaneous saturated assignments at these physical
prefixes. Their cumulative shape parameters have nonnegative increments
in $\mathfrak a$ and $\mathfrak u$, and can be selected so that
\begin{equation}
 \Delta\chi_i\le
 \frac{F(\kappa_i)-F(\kappa_{i-1})}{v},
 \qquad
 \chi_i\le F(\kappa_i)/v.
 \label{eq:chain-shape-bound}
\end{equation}
At each prefix the assignments have disjoint index ownership on each
physical block. Every retained event carries all its assignments. Their
used packet masses have the saturation exponents even after pulling the
finite chain back to the original event graph. Any fixed finite list of
additional partition tests can be imposed at the same time.
\end{lemma}

\begin{proof}
At the initial flat prefix, Lemma~\ref{lem:derivative-capture} gives
$\chi_0=0$. Thus $a=1$, $b=N^\tau$, and the packets are plates of mass
$\eqexp mN^\tau$. Cut to the outer model, keeping its full ensemble of
matrix parents with inherited weights, rather than selecting a single
parent. Apply derivative capture at the next physical prefix in this
ensemble, and continue.

The parameter calculation for one step explains the choice of cuts.
Suppose the preceding cumulative parameters are
$(\mathfrak a_{\rm old},\mathfrak u_{\rm old})$, with
$\chi_{\rm old}=\mathfrak a_{\rm old}+\mathfrak u_{\rm old}$, and put $L=1-\chi_{\rm old}$.
If the next cut in the normalized outer model has shape
$(\mathfrak a',\mathfrak u')$, its pullback has
\[
 (\mathfrak a,\mathfrak u)
   =(\mathfrak a_{\rm old},\mathfrak u_{\rm old})
      +L(\mathfrak a',\mathfrak u').
\]
Indeed spatial cell widths multiply. Physical time is unchanged, but
in the new cell units
\[
 K_{\rm new}=K N^{-\chi_{\rm old}},\qquad
 b_{\rm new}=b N^{-\mathfrak a_{\rm old}},\qquad
 \frac{K_{\rm new}}{b_{\rm new}}
   =\frac Kb N^{-\mathfrak u_{\rm old}}.
\]
This also gives the asserted increment in $\mathfrak u$; both increments
are nonnegative. If the preceding physical cut was at $\kappa_{\rm old}$,
the outer deficit up to the next physical scale $\kappa$ is
$[F(\kappa)-F(\kappa_{\rm old})]/L$, by \eqref{eq:outer-profile}.
The first selection in Lemma~\ref{lem:derivative-capture} therefore gives
\[
 \chi-\chi_{\rm old}\le [F(\kappa)-F(\kappa_{\rm old})]/v.
\]
Iteration gives \eqref{eq:chain-shape-bound}. Since $F$ is
$1$-Lipschitz and $v>0$ is fixed, adjacent changes in both shape
parameters are bounded by a constant times the mesh size.

We give the retention details for one step and then for its pullback.
Each current parent has bounded normalized matrix coordinates, the common
profile, and the inherited upper bounds. Summing support counts with the
different parents counted separately preserves the aggregate lower
multiplicity in \eqref{eq:outer-multiplicity}. The greedy argument of
Lemma~\ref{lem:derivative-capture} applies in each parent and physical
block with this common raw threshold. If it could not capture an
extensive piece of the ensemble, the uniform perturbed upper bound on
the remainder, summed over the parents, would contradict that aggregate
lower estimate. Retain a common width type among the captured packets,
and reunite the selected blocks of each index in its current full-time
parent. There is one such parent per index, so this reunion does not
duplicate its weight. Regularization preserves the profile and the packet
masses, after which the next isotropic parent transformation is defined.

All the derived blocks used at the next stage are coarser physical blocks
than those already filled. A filled time is present only because its
index occupied the earlier, finer block. Hence activity in a derived
block implies original activity in the same physical block. The pulled
back upper bounds therefore count subsets of the original active
indices, as required by Proposition~\ref{prop:local-bound}.

For the pullback, consider an earlier registered pair consisting of an
index and its occupied cut block. Before filling, this pair supplied a
common number in exponent of selected events; after filling and
rebinning, it supplies another common number. The index weight is
unchanged. An extensive selection of the later events therefore hits an
extensive weight of these pairs. Keep all the preceding selected events
of each hit pair. A later assignment is constant on that earlier pair:
its temporal block is a coarser member of the same nested dyadic grid,
and its ownership decision concerns the index on that block. Thus this
pullback preserves all later assignment memberships. Repeating the
operation restores an extensive graph of original events carrying all
assignments in the finite chain.

One must also preserve the masses of the used parts of earlier packets.
At any fixed earlier level, its assignments are disjoint on each physical
block, and every assigned index had the same event count in exponent
there. A packet's event mass is consequently a common factor times its
index mass. Include the packet partitions at every preceding level in
the deficiency deletion of Lemma~\ref{lem:regularization}. For a finite
list, choose the deletion cutoff smaller than the total retained
proportion, while still subpower. The charge from each partition is at
most its cutoff times the original total event mass. Iterating deficient
cell deletion over the finite list leaves an extensive graph and retains
each packet's lower mass exponent. Any other fixed finite list of tests
is included in this same deletion procedure. When only an upper packing
count is needed, one may also keep the whole earlier assigned index set,
which remains a valid active set for the original local bound.

Selections obtained by a derivative limit are made before the next cut:
for each fixed perturbation and tolerance perform the finite selection,
then take a subsequence attaining its limiting shape. Only finitely many
such choices occur in the present lemma. At any stage at which a pure
outer equality is needed, rather than its full ensemble, one may choose
a representative parent subsequence: the inherited upper bounds are
uniform at a fixed positive remaining length, and at least one parent
attains the aggregate lower estimate in exponent. The construction
itself continues with the ensemble. Finally, each $\kappa_i<1$ and
$\chi_i\le\kappa_i$ leave positive outer length, so no zero-length
iteration is used.
\end{proof}

In what follows, a limit over finer meshes always has the following
order. Fix finitely many scales, tests, and comparison tolerances; carry
out Lemma~\ref{lem:nested-cuts}; then take the resolution sufficiently
large to absorb their exponent errors. Increase the finite list and
refine the mesh only along a subsequent slow diagonal. Thus no estimate
requires a preassigned resolution error to be smaller than every
shrinking scale gap at once.

At a cut of scale $s$, the packet's aspect exponent is $s-\chi$.
The selected bound $\chi\le F(s)/v$ makes it at least
$s-F(s)/v$. The next packing argument shows that the corresponding
angular tolerance controls how much the packet normal can change at
that step. Along a chain the largest tolerance, hence the smallest
aspect exponent, controls the total change. This is why a running
minimum enters the inner rescaling.

\begin{lemma}[Packing and inner equality]
\label{lem:angle-packing}
For a nontrivial equality with delay $\tau$, fix $\tau<\kappa<1$ and set
\begin{equation}
 D_*(\kappa)=\min_{\tau\le s\le\kappa}
                   \bigl(s-F(s)/v\bigr).
 \label{eq:running-orientation-exponent}
\end{equation}
Then $0\le D_*(\kappa)\le\tau$. On an extensive refinement, the normals
of the assigned initial plates used by one index in a $\kappa$-block lie
within $N^{-D_*(\kappa)+o(1)}$ of one normal for that index and block.
If $\kappa-D_*(\kappa)>0$, there is an inner equality of effective
length exponent $\kappa-D_*(\kappa)$, initial delay
\begin{equation}
 \frac{\tau-D_*(\kappa)}{\kappa-D_*(\kappa)},
 \qquad\text{and deficit}\qquad
 \frac{F(\kappa)}{\kappa-D_*(\kappa)}.
 \label{eq:inner-delay-deficit}
\end{equation}
The corresponding conclusions hold at endpoint prefixes by limits from
within the required positive lengths.
\end{lemma}

\begin{proof}
\emph{Packing adjacent packets.}
Use the finite chains of Lemma~\ref{lem:nested-cuts}. Write $\delta$ for
their mesh size. Consider one smaller-block packet of widths $(a,b)$ and
the larger-block packets meeting it along assigned indices at dyadic
angle $\theta\gg a/b$. Adjacent widths differ by $N^{O(\delta)}$.
Each meeting is witnessed by an affine index belonging to both packets,
so their centers are tied to that same segment throughout the smaller
block. At a fixed time, a strip of length comparable to $b$ tilted by
$\theta$ has normal spread $O(a+\theta b)=O(\theta b)$ relative to
the smaller packet. Extrapolating the shared segment to the larger
block costs another factor $N^{O(\delta)}$. Thus all these larger
packets lie in a plank of widths
$N^{O(\delta)}(\theta b,b)$, aligned with the smaller packet. Their
assigned index sets on the larger block are disjoint. Dividing the local
upper mass of this enclosing plank by their saturated mass bounds their
number by
\[
 N^{O(\delta)+o(1)}\theta^x(b/a)^x.
\]
On the smaller block, the two thin conditions bound transverse motion
along their intersection by $O(a/\theta)$: solving the two strip
inequalities divides by the sine of their angle. Their intersection
is therefore contained in a plank of widths
$N^{O(\delta)}(a,a/\theta)$. Figure~\ref{fig:angle-packing} shows
the two width calculations. Apply the local upper bound again and divide
by the smaller packet's saturated mass. The fraction of its assigned
indices accounted for at this angle is at most
\begin{equation}
 N^{O(\delta)+o(1)}
 \left[\theta^x(b/a)^x\right]
 \frac{a^x(a/\theta)^y}{a^xb^y}
 =N^{O(\delta)+o(1)}
       \left(\frac{a}{\theta b}\right)^{y-x}.
\label{eq:adjacent-angle-packing}
\end{equation}

\begin{figure}[!htbp]
\centering
\begin{tikzpicture}[x=1.05cm,y=1.05cm,>=Stealth,font=\small]
\begin{scope}
\draw[dashed,gray] (-2.15,-.76) rectangle (2.15,.76);
\filldraw[fill=linkblue!12,draw=linkblue] (-2,-.10) rectangle (2,.10);
\begin{scope}[rotate=18]
\filldraw[fill=orange!15,draw=orange!80!black] (-2,-.10) rectangle (2,.10);
\end{scope}
\fill (0,0) circle (1.5pt);
\draw[->] (.7,0) arc[start angle=0,end angle=18,radius=.7];
\node at (1.03,.16) {$\theta$};
\draw[<->] (-2.15,-1.05)--(2.15,-1.05)
 node[midway,below] {$b$};
\draw[<->] (-2.45,-.76)--(-2.45,.76)
 node[midway,left] {$\theta b$};
\node at (0,1.15) {Enclosing widths};
\end{scope}
\begin{scope}[xshift=7cm]
\filldraw[fill=linkblue!12,draw=linkblue] (-2,-.10) rectangle (2,.10);
\begin{scope}[rotate=18]
\filldraw[fill=orange!15,draw=orange!80!black] (-2,-.10) rectangle (2,.10);
\end{scope}
\filldraw[fill=linkblue!40,draw=linkblue!80!black]
 (-.6314,-.1)--(.0158,-.1)--(.6314,.1)--(-.0158,.1)--cycle;
\draw[dashed,gray] (-.66,-.13) rectangle (.66,.13);
\draw[<->] (-.66,-.55)--(.66,-.55)
 node[midway,below] {$a/\theta$};
\draw (2.3,-.1)--(2.3,.1);
\draw (2.23,-.1)--(2.37,-.1);
\draw (2.23,.1)--(2.37,.1);
\node[right] at (2.37,0) {$a$};
\node at (0,1.15) {Intersection widths};
\end{scope}
\end{tikzpicture}
\caption{The transverse geometry in the adjacent-packet estimate,
schematically at comparable scales. A common assigned affine index ties
the packet centers together; the dot in the left panel represents its
trace. For $\theta\gg a/b$, the enclosing normal width is of order
$\theta b$, while the intersection of the two thin strips has length
of order $a/\theta$. Endpoint extrapolation and the adjacent scale
change contribute the $N^{O(\delta)}$ factors in the proof. All
dimensions in the diagram are understood up to constant factors.}
\label{fig:angle-packing}
\end{figure}


The factors omitted in the displayed quotient are changes of time
length, profile cost, and width type across one mesh step; their
logarithms are $O(\delta)$. If an extrapolated width passes the full
spatial size, cover it by the corresponding $N^{O(\delta)}$ bounded
planks. The same bound holds for event mass because the active indices
in a packet have uniform event counts on its block.

Since $y-x>0$, angles larger by a fixed positive exponent than $a/b$ can
be discarded with a power-small loss, after making $\delta$ sufficiently
small. Include the earlier packet partitions in the subsequent cleaning.
At cut $i$ the aspect exponent is $\kappa_i-\chi_i$, which is at least
$\kappa_i-F(\kappa_i)/v$. Summing the angular errors along the finite
chain, rather than adding their exponents, bounds the total angle by
$N^{-D_*(\kappa)+O(\delta)+o(1)}$. The final packet has one owner per
index on its physical block, so all retained initial plates of that
index on the block compare to the same final normal. Sending the mesh
to zero in the specified order proves the orientation assertion.

\paragraph{The inner equality.}
The reverse selection in Lemma~\ref{lem:derivative-capture} has
$\chi\ge F(s)/v$, while every admissible shape has $\chi\le s$.
Therefore $s-F(s)/v\ge0$. Its value at $s=\tau$ is $\tau$, proving the
bounds on $D_*$. Apply Lemma~\ref{lem:anisotropic} over the
$N^\kappa$-block with angular scale $\theta=N^{-D_*}$. Its fine assigned
plates are the flat-prefix plates of length $N^\tau$, and
$D_*\le\tau$ ensures $\theta\ge N^{-\tau}$. The effective resolution
is $N^{\kappa-D_*}$ and the raw lower multiplicity is
$mN^{D_*-o(1)}$. By the local upper bound its full-time plank parameter
is at most
\[
 mN^{-h\kappa-\Pi(\kappa)+yD_*}.
\]
Only an initial flat log interval of length $D_*$ is dropped. Multiplying
this parameter by the critical factor
$N^{h(\kappa-D_*)+\Pi(\kappa)}$ gives
$mN^{(y-h)D_*}=mN^{D_*}$, exactly the lower exponent. The normalized
profile thus gives \eqref{eq:inner-delay-deficit}. As for outer parents,
one can extract a representative equality from the aggregate chart
ensemble whenever this statement is used as a fact about equality
profiles.
\end{proof}

\subsection{A positive minimum delay and the canonical profile}

\begin{proposition}[Canonical equality profile]
\label{prop:canonical}
There is an exact equality whose profile is
\begin{equation}
 F(\kappa)=\beta(\kappa-\tau)_+,
 \qquad 0<\tau<1,\qquad
 \ell=1-\tau,\qquad \beta=v/\ell\le1.
 \label{eq:canonical-profile}
\end{equation}
Its deficit is $v$, and $\tau$ is the least initial delay among all
equalities with deficit $v$.
\end{proposition}

\begin{proof}
\emph{Attaining the maximal deficit.}
First, equalities with deficit $v$ exist. Start with any positive-deficit
equality and put $g(s)=s-F(s)/v$. If the running minimum of $g$ from
$\tau$ remained equal to $\tau$ throughout its nontrivial part,
Lemma~\ref{lem:angle-packing}, applied at a prefix with positive deficit,
would produce a nontrivial equality with zero initial delay. This
contradicts Lemma~\ref{lem:flat-delay}. There is therefore a new minimum
$D<\tau$, attained at a prefix $\kappa$ with
$g(\kappa)=D$. At this prefix the inner deficit is
\[
 \frac{F(\kappa)}{\kappa-D}=v.
\]
Let $\mathcal E_v$ denote the nonempty class of equality profiles with
deficit $v$, and set
\[
 \tau_* =\inf_{F\in\mathcal E_v}
             \sup\{r:F=0\text{ on }[0,r]\}.
\]

\paragraph{A positive lower bound for the delay.}
We claim that $\tau_*>0$. Suppose otherwise, and choose
$F_j\in\mathcal E_v$ with initial delays $\tau_j\to0$. Equality
compactness, as in Lemma~\ref{lem:hybrid-equality}, gives a subsequence
converging uniformly to an equality $F_\infty$ with deficit $v>0$.
The cost is lower semicontinuous; a strict downward jump would give a
violation of the critical inequality, so this limiting configuration is
still an equality. By Lemma~\ref{lem:flat-delay}, $F_\infty$ has a flat
interval of positive length.

The small-prefix multilinear argument underlying
Lemma~\ref{lem:prefix-extension} has slack depending only on the fixed
critical parameters. In particular, its allowed prefix length can be
chosen independently of the equality profile: the estimate used there is
\begin{equation}
 \frac{k^2m}{n}
 \ge N^{h+q+\Pi(1)-2\alpha-o(1)}
 \ge N^{h+(p-2)\alpha-o(1)}
 \ge N^{h-o(1)}.
 \label{eq:uniform-prefix-slack}
\end{equation}
Indeed $n\leexp\Delta_hN^d$,
$m\eqexp\Delta_hN^{h+\Pi(1)}$, and
$\Pi(1)\ge p\alpha-q$. Choose fixed $0<a_0<b_0$ inside both that
allowed prefix range and the flat interval of $F_\infty$.

For a sufficiently late $j$,
\[
 \tau_j<a_0,\qquad
 F_j(b_0)<\frac h4(b_0-a_0).
\]
Choose a minimizer $\kappa_j$ of $F_j(u)-(h/4)u$ on $[0,b_0]$.
At $b_0$ this function is strictly less than $-ha_0/4$, whereas for
$u\le a_0$ it is at least $-ha_0/4$. Hence $\kappa_j>a_0>\tau_j$,
so $\delta_j=F_j(\kappa_j)>0$. The minimizing property gives
\begin{equation}
 F_j(\kappa_j)-F_j(u)
 \le\frac h4(\kappa_j-u),\qquad 0\le u\le\kappa_j.
 \label{eq:late-prefix-record}
\end{equation}
Fix this particular $j$, its prefix $\kappa_j$, and its positive
$\delta_j$ before passing to the realizing resolution sequence.

Apply Lemma~\ref{lem:prefix-extension} to this fixed equality and
prefix, with $s=1-h/4$ and $\gamma=0$. Its uniform small-prefix
range includes $\kappa_j$, by the choice of $b_0$. The record inequality
\eqref{eq:late-prefix-record} is precisely its temporal hypothesis;
the other hypothesis is simply $F_j\ge0$. Its angular exponent is $h$,
so
\[
 t=1+3h/4,\qquad
 \mu=\min\{1,(s+t)/2,t\}=1.
\]
The lemma gives $\Pi(\kappa_j)\le F_j(\kappa_j)=\delta_j$.
But $\Pi(\kappa_j)\ge(p-1/2)\delta_j>\delta_j$, since
$\delta_j>0$ and $p>2$. This is a contradiction. The positive gap
is fixed after choosing $j$ and before taking its realizing resolution
to infinity; no uniform lower bound for the numbers $\delta_j$ is
needed. We have proved $\tau_*>0$.

\paragraph{Rigidity from the inner and outer delays.}
We now determine a profile attaining this infimum. For any
$F\in\mathcal E_v$ with initial delay $\tau$, $g(1)=0$. Given
$D'\in(0,\tau)$, let $s$ be the first passage below $D'$ after $\tau$:
\[
 s=\inf\{r>\tau:g(r)<D'\}.
\]
Continuity gives $\tau<s<1$, $g(s)=D'$, and
$\min_{\tau\le r\le s}g(r)=D'$. There are values below $D'$ arbitrarily
soon to the right of $s$. The inner equality from
Lemma~\ref{lem:angle-packing} has deficit $v$, so its delay is at least
$\tau_*$:
\begin{equation}
 \tau-D'\ge\tau_*(s-D').
 \label{eq:canonical-inner-passage}
\end{equation}

Any saturated outer cut at $s$ has $\chi\le s<1$ and deficit
$(v-F(s))/(1-\chi)$. By Lemma~\ref{lem:derivative-capture} this is at
most $v$, whence $\chi\le F(s)/v$. The reverse selection of that lemma
therefore attains $\chi=F(s)/v=s-D'$, and its outer deficit is exactly
$v$. Moreover $F$ increases arbitrarily soon to the right of $s$:
$g(r)<g(s)$ implies $F(r)-F(s)>v(r-s)$. Formula
\eqref{eq:outer-profile} consequently gives the exact outer delay
$D'/(1-s+D')$, not merely a lower bound for it. Its membership in
$\mathcal E_v$ yields
\begin{equation}
 D'\ge\tau_*(1-s+D').
 \label{eq:canonical-outer-passage}
\end{equation}

The two inequalities compare complementary rescalings of the same
profile, both with deficit $v$. For a sequence approaching the least
delay, they will force the same first-passage location.

Choose $F_j\in\mathcal E_v$ with $\tau_j\to\tau_*$. The Lipschitz
bound gives $v\le1-\tau_j$, so $\tau_*\le1-v<1$.
For each fixed $D'\in(0,\tau_*)$, the two passage inequalities give
\[
 1+D'-D'/\tau_*\le s_j
 \le D'+(\tau_j-D')/\tau_*.
\]
Thus
\[
 s_j\longrightarrow 1-\frac{1-\tau_*}{\tau_*}D',
 \qquad F_j(s_j)=v(s_j-D').
\]
Pass to a compact equality limit $F$. The rational values of $D'$ in
$(0,\tau_*)$ suffice; continuity then determines the profile at every
point above $\tau_*$. Below $\tau_*$ all these profiles vanish. We obtain
\[
 F(\kappa)=\frac{v}{1-\tau_*}(\kappa-\tau_*)_+.
\]
Its delay is exactly $\tau_*$ and its slope is at most one by the
Lipschitz bound. This proves \eqref{eq:canonical-profile}.
\end{proof}

Figure~\ref{fig:canonical-profile} shows the canonical profile and the
aspect exponent that controls its angular scale.

\begin{figure}[H]
\centering
\begin{tikzpicture}[x=4.6cm,y=3.0cm,>=Stealth,font=\small]
\begin{scope}
\draw[->] (0,0)--(1.1,0) node[right] {$\kappa$};
\draw[->] (0,0)--(0,1.08) node[above] {$F(\kappa)$};
\draw[very thick,linkblue] (0,0)--(.36,0)--(1,.72);
\draw[dashed,gray] (1,0)--(1,.72)--(0,.72);
\node[left] at (0,.72) {$v$};
\node[below] at (.36,0) {$\tau$};
\node[below] at (1,0) {$1$};
\node[above left] at (.8,.5) {slope $\beta=v/\ell$};
\node[below=16pt] at (.55,0) {Canonical deficit profile};
\end{scope}
\begin{scope}[xshift=7.0cm]
\draw[->] (0,0)--(1.1,0) node[right] {$\kappa$};
\draw[->] (0,0)--(0,1.08) node[above] {$\kappa-F(\kappa)/v$};
\draw[very thick,linkblue] (0,0)--(.36,.55)--(1,0);
\draw[dashed,gray] (.36,0)--(.36,.55)--(0,.55);
\node[left] at (0,.55) {$\tau$};
\node[below] at (.36,0) {$\tau$};
\node[below] at (1,0) {$1$};
\node[above] at (.75,.3) {$\tau(1-\kappa)/\ell$};
\node[below=16pt] at (.55,0) {Remaining orientation exponent};
\end{scope}
\end{tikzpicture}
\caption{The canonical profile in Proposition~\ref{prop:canonical}.
After the flat interval, its slope is constant. For $\kappa\ge\tau$, the associated quantity
$\kappa-F(\kappa)/v$ decreases linearly from $\tau$ to zero;
at the stationary prefix $\kappa=1-\ell c$ it equals $\tau c$.
This identity gives the angular scale $\Theta_c=N^{-\tau c}$.
The vertical scales in the two panels are independent.}
\label{fig:canonical-profile}
\end{figure}


\subsection{Tracking a stationary shape}

Fix the canonical equality from Proposition~\ref{prop:canonical}. For
the remainder of the proof, $\ell$ always denotes $1-\tau$; the
logarithmic loss of an outer cut remains $\chi$.

\begin{lemma}[Canonical shape tracking]
\label{lem:stationary-shape}
For the canonical profile there are compatible finite chains, and slow
mesh limits of them, with cumulative shape parameters
\begin{equation}
 \kappa=\tau+\ell b',\qquad
 \chi=b',\qquad \mathfrak a=wb',\qquad
 w=h_z+\ell\mathcal D_q(q,\beta)\in(0,1],
 \quad 0\le b'<1.
 \label{eq:stationary-shape}
\end{equation}
The equalities in the parameters allow errors tending to zero with the
mesh. Packet mass, ownership, and orientation conclusions of
Lemmas~\ref{lem:nested-cuts} and~\ref{lem:angle-packing} remain valid.
\end{lemma}

\begin{proof}
An outer cut of a deficit-$v$ equality again has deficit at most $v$.
Applied to the canonical profile, this forces $\chi\le b'$. The reverse
derivative selection attains $\chi=b'$. Formula
\eqref{eq:outer-profile} then shows that the outer profile is the same
canonical profile: its flat delay is
\[
 \frac{\tau+\ell b'-b'}{1-b'}=\tau,
\]
and its nonflat slope is still $\beta$. The same statement applies
incrementally in every current outer model, in the original logarithmic
units.

For one mesh increment $\Delta b'$, the two coefficients in
\eqref{eq:derivative-capture}, written in those units, are
\[
 v(\Delta b'-\Delta\chi),\qquad
 h_z\Delta\chi+\ell\mathcal D_q(q,\beta)\Delta b'
                  -\Delta\mathfrak a.
\]
Every available saturation has $\Delta\chi\le\Delta b'$ by the outer
deficit bound. Take either sign of $\dot z$ and take $\dot p<0$ with
an arbitrarily large ratio $|\dot p|/|\dot z|$. First fix that ratio and
perform derivative capture. Along a subsequence as the ratio increases,
the nonnegative quantity $\Delta b'-\Delta\chi$ must tend to zero;
otherwise its negative $\dot p$ contribution cannot be offset by the
bounded second coefficient. The two signs of $\dot z$ therefore provide
choices with $\Delta\chi=\Delta b'$ and, respectively,
\[
 \Delta\mathfrak a\le w\Delta b',
 \qquad\text{or}\qquad
 \Delta\mathfrak a\ge w\Delta b'.
\]
Since $0\le\Delta\mathfrak a\le\Delta\chi$, this implies
$0\le w\le1$. Its defining expression and $h_z>0$ give $w>0$.

It is unnecessary to realize an exact intermediate shape on a single
step. If the running value of $\mathfrak a$ is above $wb'$, choose the
next increment from the lower side; if below, choose from the upper
side. The error can then overshoot zero by at most the mesh size, because
both $\Delta\mathfrak a$ and $w\Delta b'$ lie in
$[0,\Delta b']$. This tracks the desired total shape within one mesh
increment. The current outer profile stays canonical after each choice,
so both sides are available at every step. All perturbation limits are
resolved on a fixed finite chain before taking its mesh diagonal, and
Lemma~\ref{lem:nested-cuts} preserves the required earlier assignments.
\end{proof}

We now express the tracked packets in physical scales. Reverse the cut
parameter by $c=1-b'$ and put
\[
 \rho=1-w,\qquad B_*=h+\ell P_{p,q}(\beta),\qquad
 R_c=N^{-\ell c},\qquad\Theta_c=N^{-\tau c}.
\]
The corresponding logarithmic scale and cumulative shape are
\[
 \kappa=1-\ell c,\qquad \chi=1-c,\qquad
 \mathfrak a=w(1-c),\qquad \mathfrak u=\rho(1-c).
\]
For $0\le c\le1$, use nested dyadic time partitions whose blocks $I_c$
have physical length $R_c$, up to dyadic rounding. For
$0<c\le1$, write $B_c$ for its assigned packet. Its thin and thick
spatial widths in original cell units are
\begin{equation}
 a=N^{w(1-c)},\qquad
 b=N^{w+(\tau-w)c}.
 \label{eq:stationary-spatial-widths}
\end{equation}
At $c=1$ these are the initial flat-scale plates. As $c$ decreases,
the time blocks grow and the packets record the successive outer cuts.

For the probability statements, replace the finer events of each
occupied index--$I_1$ pair by one atom. More precisely, let
$\mathcal J_i$ be the occupied $I_1$ blocks on index $i$ and set
\[
 \Omega_* =\{(i,J):J\in\mathcal J_i\},\qquad
 Z_* =\sum_i\omega_i\#\mathcal J_i,\qquad
 \Prob\{(i,J)\}=\frac{\omega_i}{Z_*}.
\]
The atom $(i,J)$ is placed at the center $t_J$ of $J$ and carries the
spatial normal of the initial plate assigned to that index on $J$.
We restrict these normals to a fixed direction chart. After a fixed
orthonormal spatial change of coordinates, write their covectors as
$(1,\vartheta)$ with $\vartheta$ bounded. In the displayed law,
$\mathcal J_i$ and $Z_*$ refer to the retained pairs. Denote the atom's
time and normal coordinates by $e=(t,\vartheta)$, while retaining $i$
as a separate random variable. Nested dyadic quantization of $t$ and
$\vartheta$ at widths $R_c$ and $\Theta_c$ gives labels $e_c$.
The original uncollapsed equality supplies $m$, $\Delta_h$, and all
active-index plank bounds; this collapsed law supplies event probabilities.

In an event frame define the four matrix forms
\begin{equation}
 X=(1,\vartheta)\cdot M(t),\qquad Y=M(t)_2,
 \qquad U=(1,\vartheta)\cdot u,\qquad V=u_2.
 \label{eq:stationary-matrix-forms}
\end{equation}
Thus $X,U$ are normal position and velocity, and $Y,V$ are the fixed
complementary coordinates. This chart is generally oblique. Dividing
\eqref{eq:stationary-spatial-widths} by $N$ gives the spatial widths;
dividing again by the time length $R_c$ gives the velocity widths.
We write these four scales as
\begin{equation}
\begin{aligned}
 L_X(c)&=N^{-\rho-wc},&
 L_Y(c)&=L_X(c)\Theta_c^{-1},\\
 L_U(c)&=L_X(c)R_c^{-1},&
 L_V(c)&=L_X(c)(R_c\Theta_c)^{-1}.
\end{aligned}
\label{eq:estimate-widths}
\end{equation}
The proposition verifies that these widths remain valid in every used
event frame, and records the masses and branching of the same packets.

\begin{proposition}[Stationary data]
\label{prop:stationary-data}
There is a realizing sequence of canonical equalities and compatible
packet assignments for which the preceding collapsed law, event labels,
and matrix forms have the following properties. All local upper bounds
and original index weights are retained.
\begin{enumerate}[label=\textup{(\roman*)}]
\item
Along an index, occupied time blocks have regular branching exponent
$s_0=1-\beta\ge0$. More precisely, for fixed $c<b\le1$, an occupied
$I_c$ contains
\begin{equation}
 N^{\ell(1-\beta)(b-c)+o(1)}
 \label{eq:stationary-time-branching}
\end{equation}
occupied $I_b$ blocks. The event normals on that index in $I_c$ lie in a
range of size $\Theta_cN^{o(1)}$.

\item
For $0<c\le1$ there are assigned packets $B_c$ on the occupied $I_c$
blocks, with raw active index mass at least
\begin{equation}
 n_c=\Delta_h
 N^{(B_*+\tau)c+dw(1-c)}
 \label{eq:stationary-packet-mass}
\end{equation}
in exponent. On each fixed time block the assigned index sets are
disjoint, and each retained index has one packet owner there. The used
normals of a packet differ from its spatial normal by at most
$\Theta_cN^{o(1)}$.

\item
In every used event frame, the widths of $B_c$ in the four forms
\eqref{eq:stationary-matrix-forms} are bounded by
$L_X(c),L_Y(c),L_U(c),L_V(c)$ up to subpower expansion. Equivalently,
their negative log exponents are
\begin{equation}
 \rho+wc,\qquad \rho+(w-\tau)c,
 \qquad\rho+(w-\ell)c,\qquad\rho+(w-1)c,
 \label{eq:stationary-matrix-widths}
\end{equation}

\item
The data and packet lower bounds coexist on any prescribed finite list
of scales and tests, and hence on countable lists by slow diagonals. They
persist, in exponent, on typical nondeficient parts after extensive
regular refinements. All packet memberships are memberships of original
active indices on their physical intervals. In a chart registered at
$I_c$, the packet at a finer $I_b$ is determined by the index, chart, and
$I_b$; a matrix cell in that chart is determined by the index and chart.
\end{enumerate}
\end{proposition}

\begin{proof}
Choose the finite chains from Lemma~\ref{lem:stationary-shape} and
use the scales defined above.
On the canonical profile the running aspect exponent is decreasing
above $\tau$, and at $\kappa=1-\ell c$ its value is
\[
 \kappa-F(\kappa)/v=\tau c.
\]
The proof of Lemma~\ref{lem:angle-packing}, which applies unchanged to
the tracked finite chains, shows that every initial plate used by this
packet has its normal within $\Theta_cN^{o(1)}$ of the packet normal.
Since a packet is assigned only once to each index on its block, this
also gives the range assertion for all event normals on that index and
block. The canonical profile gives the number of occupied subblocks:
subtracting its deficits between the prefixes $1-\ell c$ and
$1-\ell b$ yields \eqref{eq:stationary-time-branching}.

We now verify the width bounds in the event frames.
The widths in \eqref{eq:estimate-widths} are first computed in the
packet's own normal frame. Changing to one of its used event normals
adds at most $\Theta_cN^{o(1)}$ times the transverse coordinate; the
ratios of transverse to normal widths are exactly $\Theta_c^{-1}$.
Changing the time center within $I_c$ adds at most $R_c$ times the
velocity coordinate. Both changes preserve the four widths up to
subpower factors.

For this profile,
\[
 \Pi(\kappa)=\ell(1-c)P_{p,q}(\beta),\qquad
 m\eqexp\Delta_hN^{B_*}.
\]
Substituting \eqref{eq:stationary-spatial-widths} in
\eqref{eq:stationary-local-threshold}, the exponent of the packet mass
relative to $\Delta_h$ is
\begin{align*}
 &B_*-\ell(1-c)P_{p,q}(\beta)-h(1-\ell c)
      +xw(1-c)+y\bigl(w+(\tau-w)c\bigr)\\
 &\hspace{12mm}=dw(1-c)
        +c\bigl(\ell P_{p,q}(\beta)+h\ell+y\tau\bigr)\\
 &\hspace{12mm}=dw(1-c)+(B_*+\tau)c,
\end{align*}
using $\ell+\tau=1$ and $y=1+h$. This proves
\eqref{eq:stationary-packet-mass}.

We next verify the passage to the collapsed law. Before collapse,
the number of finer events in every occupied index--$I_1$ pair has a
common exponent. Thus an extensive selection of atoms of weight
$\omega_i$ pulls back to an extensive selection of the uncollapsed
graph. Restricting the normal to one of finitely many direction charts
also retains an extensive piece, and normalized covectors and unit
normals are comparable on that chart. The orientation and time ranges
proved above imply that one index and $I_c$ use only subpower many
$e_c$ bins. Whenever the chart restriction or these labels split a
packet, retain nondeficient occupied pieces. The packet lower bounds
therefore survive for the law specified before the proposition.

For completeness, scale interpolation does not require new unrelated
assignments. At a finite mesh, use the packets already constructed.
A nearby intermediate interval is assigned the packet of its nearest
coarser mesh interval, restricted to the indices active in the smaller
interval. Extrapolation, width changes, and subdivision have
$N^{O(\text{mesh})}$ factors. Regular time counts and deficiency deletion
preserve the packet lower masses on typical occupied pieces. Along the
slow diagonal these factors are subpower. Alternatively, include any
specified intermediate scales in the next finite mesh. Thus a new
finite or countable list of scale tests can be accommodated by further
refinement, with the same limiting data.

Finally include, in the finite cleaning list, the assignment partitions
from every level in use. Lemma~\ref{lem:nested-cuts} then gives all the
simultaneous memberships, the original active-index interpretation, and
stability under extensive regular refinement in (iv). On a fixed
physical $I_b$, ownership was a decision for an index on that block;
after registration in a chart it is therefore constant on the tuple
specified in (iv). A cell of that chart's fixed matrix partition is
likewise a function only of the index and chart. The construction needs
outer steps only for $b'<1$, or $c>0$; it makes no cut at zero remaining
length. The endpoint $c=1$ is the initial assignment and has already been
included.
\end{proof}

We have proved $0<\tau<1$, $0<\ell<1$, $v>0$, and
$\beta=v/\ell>v$. The value $w=1$ is allowed, so that $\rho$ may
vanish. At this stage $\beta=1$ has not been excluded; the angular
estimate in the next section will show that $s_0=1-\beta$ is positive.
